% ===========================================================
% Chapter 15: Reasoning, Tool Use, and Reinforcement Finetuning
% Part IV: Frontiers
% ===========================================================

\chapter[Рассуждения и дообучение с подкреплением]{Рассуждения, инструменты и дообучение с подкреплением}
\label{ch:reasoning}

\begin{quote}
\itshape
Если интеллект~--- это торт, то основная его часть~--- это обучение без учителя,
глазурь~--- это обучение с учителем, а вишенка на торте~---
обучение с подкреплением.
\par\medskip
\upshape\raggedleft
---\,Yann LeCun, NeurIPS 2016
\end{quote}

\bigskip

На протяжении большей части истории, описанной в этой книге, обучение с подкреплением
занимало любопытное место в конвейере обучения языковых моделей:
теоретически центральное, практически маргинальное. Сигнал вознаграждения, который
управлял RLHF\index{RLHF} (Глава~\ref{ch:practical-rlhf}), сам по себе являлся
обученной нейронной сетью,\,---\,несовершенным приближением для
несовершенно понятой цели,\,---\,и каждый практик знал, что слишком сильное
смещение стратегии в сторону такой модели вознаграждения приведёт к
взлому вознаграждения, а не к подлинному улучшению. Результатом стал
осторожный, жёстко ограниченный по KL режим обучения, в котором RL вносил
лоск, но не мощь.

Эта картина резко изменилась между 2024 и 2025 годами. Катализатором послужило
осознание того, что определённые области допускают \emph{верифицируемые} вознаграждения:
математический ответ либо верен, либо нет; фрагмент кода
либо проходит набор тестов, либо нет. Когда вознаграждения можно
вычислить с помощью детерминированной проверочной функции, а не выводить из
нейронной модели предпочтений, опасность взлома вознаграждения исчезает,
и становится безопасным применять значительно более агрессивное RL на значительно более
длительный срок. Результат\,---\,\emph{Обучение с подкреплением с верифицируемыми
вознаграждениями} (RLVR)\index{RLVR}\,---\,произвёл качественный скачок в
способности к рассуждению: модели теперь тратят тысячи токенов на обдумывание
перед ответом на сложный вопрос, могут писать и отлаживать сложные
программы и достигают экспертного уровня на олимпиадной математике.

Эта глава организована для постепенного формирования понимания. Мы начинаем
с формализации цели RLVR и сравнения её с
целью RLHF из Части~III
(Раздел~\ref{sec:rlvr-ch15}). Затем мы рассматриваем рассуждения типа chain-of-thought\index{chain-of-thought}
как технику формирования промптов и как обучаемое
поведение (Раздел~\ref{sec:cot}). Раздел~\ref{sec:training-reasoning}
изучает рецепты обучения, используемые на практике, с DeepSeek R1 в качестве
разобранного примера. Раздел~\ref{sec:inference-scaling} формализует
масштабирование на этапе вывода: примечательное явление, при котором генерация большего числа
токенов во время тестирования приводит к измеримому росту точности.
Разделы~\ref{sec:tool-use} и~\ref{sec:agentic-rl} расширяют
рамки на \emph{использование инструментов}\index{tool use} и \emph{агентное
RL}\index{agentic RL}, где модель должна взаимодействовать с
внешними средами в течение нескольких шагов. Раздел~\ref{sec:practical-reasoning}
обозревает практические алгоритмические и инфраструктурные решения, а
Раздел~\ref{sec:reasoning-summary} завершает открытыми вопросами.

% ===========================================================
\section{От RLHF к обучению с подкреплением с верифицируемыми вознаграждениями}
\label{sec:rlvr-ch15}
% ===========================================================

\subsection{Инсайт верификации}

Напомним из Главы~\ref{ch:practical-rlhf}, что RLHF максимизирует
\begin{equation}
\label{eq:rlhf-objective-ch15}
J_{\mathrm{RLHF}}(\theta)
  = \E_{\pi_\theta}\bigl[r_\phi(x, y)\bigr]
  - \beta\,\KL{\pi_\theta}{\pi_{\mathrm{ref}}},
\end{equation}
где $r_\phi$~--- это \emph{обученная} модель вознаграждения, обученная на парах предпочтений
людей, а $\beta > 0$~--- коэффициент штрафа KL, который
предотвращает слишком сильное отклонение стратегии от референсной модели.
Обученная модель вознаграждения порождает две взаимоусиливающиеся проблемы. Во-первых, это
приближение: даже хорошо обученная модель вознаграждения будет присваивать
неверные оценки некоторой доле ответов, и достаточно сильная стратегия может
обнаружить и использовать эти ошибки\,---\,явление,
называемое \emph{сверхоптимизацией вознаграждения}\index{reward overoptimisation}
или законом Гудхарта в RL (Глава~\ref{ch:reward-models}). Во-вторых,
обучение модели вознаграждения требует дорогостоящей разметки людьми пар
предпочтений, и полученная модель наследует все предвзятости и
несогласованности этих разметок.

\index{RLVR!definition}
\emph{Обучение с подкреплением с верифицируемыми вознаграждениями} (RLVR) обходит
обе проблемы, заменяя обученную модель вознаграждения детерминированной
\textbf{функцией верификации}\index{verification function}:
\begin{equation}
\label{eq:rlvr-reward-ch15}
r(x, y) = \begin{cases}
  1 & \mbox{если } \textsc{Verify}(x, y) = \textsc{True}, \\
  0 & \mbox{иначе.}
\end{cases}
\end{equation}
Функция $\textsc{Verify}$ является предметно-специфической и полностью
символической: для математики она извлекает окончательный ответ из
ответа модели (используя регулярное выражение для нахождения содержимого в
нотации \texttt{\textbackslash boxed\{\}} или после разделителя, такого
как \texttt{The answer is:}) и сравнивает его с правильным
решением; для кода она выполняет ответ на наборе тестов и
возвращает \textsc{True}, если все проверки проходят.

\begin{definition}[Цель RLVR]
\label{def:rlvr-ch15}
\index{RLVR!objective}
Пусть $\mathcal{D}$~--- это набор данных пар (промпт, правильный ответ)
$(x, a^*)$. Цель RLVR:
\begin{equation}
\label{eq:rlvr-objective}
J_{\mathrm{RLVR}}(\theta)
  = \E_{(x,a^*)\sim\mathcal{D},\,y\sim\pi_\theta(\cdot|x)}
    \bigl[\ind\{\textsc{Verify}(y, a^*)\}\bigr]
  - \beta\,\KL{\pi_\theta}{\pi_{\mathrm{ref}}},
\end{equation}
где $\ind\{\cdot\}$~--- индикаторная функция, а член KL может быть
уменьшен или полностью исключён (см. Раздел~\ref{sec:training-reasoning}).
\end{definition}

Поскольку $r$ принимает значения в $\{0, 1\}$ и вычисляется
детерминированной символической функцией, она не может быть сверхоптимизирована
так же, как нейронная модель вознаграждения. Модель, научившаяся давать
верные ответы, \emph{действительно} стала лучше; нет прокси-разрыва для
использования.

\subsection{Верифицируемые области}
\label{subsec:verifiable-domains}

Мощь RLVR зависит от наличия богатого и разнообразного набора задач
с эффективно верифицируемыми ответами. В настоящее время доминируют три области.

\begin{description}[leftmargin=!,labelwidth=3.5cm]
\item[Математика] Олимпиадные задачи из наборов данных, таких как
  MATH и AIME, допускают единственные числовые или символические ответы, которые можно
  проверить точно. Модели предлагается заключить окончательный
  ответ в окружение \texttt{\textbackslash boxed\{\}}, которое регулярное
  выражение извлекает и сравнивает с эталоном.

\item[Генерация кода] Задачи с LeetCode, HumanEval или MBPP
  поставляются с модульными тестами. Решение модели выполняется в изолированной среде,
  и вознаграждение равно $1$, если все тесты проходят (или доле прошедших тестов
  для вариантов с частичным зачётом).

\item[Наука и логика] Формальное доказательство теорем (например, в Lean 4) и
  структурированное научное вопросно-ответное взаимодействие с уникальными фактическими ответами
  являются естественными расширениями. Любая область, где можно написать проверочную функцию,
  в принципе доступна для RLVR.
\end{description}

\subsection{Формальное сравнение с целью RLHF}
\label{subsec:rlvr-vs-rlhf}

Таблица~\ref{tab:rlhf-vs-rlvr} сравнивает два подхода по
параметрам, наиболее значимым для практиков.

\begin{table}[ht]
\centering
\small
\caption{RLHF против RLVR: формальное сравнение.}
\label{tab:rlhf-vs-rlvr}
\begin{tabular}{@{}lcc@{}}
\toprule
\textbf{Свойство} & \textbf{RLHF} & \textbf{RLVR} \\
\midrule
Источник вознаграждения & Обученная нейронная $r_\phi$ & Символическая $\textsc{Verify}$ \\
Диапазон вознаграждения & $\R$ (непрерывный) & $\{0, 1\}$ (бинарный) \\
Разметка людьми & Требуется для $r_\phi$ & Требуется для пар $(x, a^*)$ \\
Риск сверхоптимизации & Высокий & Пренебрежимо мал \\
Необходимость штрафа KL & Критическая & Необязательна \\
Применимые области & Общие & Верифицируемые области \\
Стоимость модели вознаграждения & $O(\mbox{длина ответа})$ & $O(\mbox{проверочная функция})$ \\
\bottomrule
\end{tabular}
\end{table}

Ключевое структурное отличие состоит в том, что RLVR преобразует задачу RL в такую,
где вознаграждение \emph{не является} нейронной сетью со своими градиентами, а
является оракулом только для чтения. Это делает ландшафт оптимизации
чище: нет внутреннего цикла обучения модели вознаграждения, который нужно чередовать
с обновлениями стратегии, нет проблемы сдвига распределения для модели
вознаграждения и нет риска того, что стратегия «найдёт пробелы» в обученном
приближении.

\begin{figure}[ht]
\centering
\begin{tikzpicture}[
  node distance=1.8cm and 2.2cm,
  box/.style={rectangle, rounded corners=4pt, draw, thick,
              minimum width=2.4cm, minimum height=0.9cm,
              font=\small, align=center},
  verify/.style={diamond, draw, thick, aspect=2,
                 font=\small\itshape, align=center,
                 inner sep=2pt},
  arrow/.style={-{Stealth[length=7pt]}, thick},
  darrow/.style={{Stealth[length=7pt]}-{Stealth[length=7pt]}, thick, dashed}
]
% Узлы
\node[box, fill=blue!10]  (prompt)    {Промпт $x$};
\node[box, fill=green!10, right=of prompt]  (generate)  {Генерация\\$y \sim \pi_\theta(\cdot|x)$};
\node[verify, fill=orange!15, right=of generate] (verify) {Верификация\\$(y, a^*)$};
\node[box, fill=red!10, below=1.4cm of verify]   (reward)   {Вознаграждение\\$r \in \{0,1\}$};
\node[box, fill=purple!10, left=2.8cm of reward] (update)   {Обновление стратегии\\$\nabla_\theta J_{\mathrm{RLVR}}$};

% Стрелки
\draw[arrow] (prompt)   -- (generate);
\draw[arrow] (generate) -- (verify);
\draw[arrow] (verify)   -- node[right]{\small да/нет} (reward);
\draw[arrow] (reward)   -- (update);
\draw[arrow] (update)   -- ++(0,2.3) -| (prompt)
    node[midway, above]{\small повторить по $\mathcal{D}$};

% Метка эталонного ответа
\node[box, fill=gray!10, above=1.2cm of verify] (gt) {Эталон $a^*$};
\draw[arrow, dashed] (gt) -- (verify);
\end{tikzpicture}
\caption{Цикл обучения RLVR. Промпт $x$ из набора данных подаётся
  текущей стратегии $\pi_\theta$, которая генерирует ответ $y$.
  Детерминированная функция верификации сравнивает $y$ с
  эталонным ответом $a^*$ и возвращает бинарное вознаграждение. Затем
  вычисляются градиенты стратегии, и цикл повторяется. Обученная модель
  вознаграждения не нужна.}
\label{fig:rlvr-loop}
\end{figure}

% ===========================================================
\section{Рассуждения типа chain-of-thought}
\label{sec:cot}
% ===========================================================

\subsection{Основная идея}

Стандартная авторегрессионная генерация отображает промпт $x$ непосредственно в
ответ $y$. Генерация типа chain-of-thought\index{chain-of-thought} (CoT)
вставляет \emph{след рассуждений}\index{reasoning trace} $z$ между
промптом и финальным ответом:
\[
x \;\longrightarrow\; z \;\longrightarrow\; a,
\]
где $z$~--- последовательность промежуточных шагов рассуждения, а $a$~--- финальный
ответ, извлечённый из конца генерации. Полная
генерация $y = (z, a)$ производится той же авторегрессионной моделью;
$z$ просто не показывается пользователю во многих развёртываниях.

\begin{example}[Chain-of-thought на олимпиадной задаче]
\label{ex:cot-math}
Рассмотрим промпт: \emph{«Найдите все целочисленные решения
$x^2 - 5x + 6 = 0$.»} Без chain-of-thought модель может выдать
правильный ответ $\{2, 3\}$ напрямую (если запомнила его) или
тихо провалиться. С chain-of-thought генерация выглядит так:

\begin{quote}
\texttt{<thinking>}\\
\texttt{Мне нужно разложить $x^2 - 5x + 6$ на множители. Ищу два числа,}\\
\texttt{произведение которых равно $6$, а сумма~--- $-5$. Это $-2$ и $-3$.}\\
\texttt{Значит, $x^2 - 5x + 6 = (x-2)(x-3)$. Приравниваем каждый множитель к нулю:}\\
\texttt{$x - 2 = 0 \Rightarrow x = 2$; $x - 3 = 0 \Rightarrow x = 3$.}\\
\texttt{</thinking>}\\
\texttt{Целочисленные решения: $x = 2$ и $x = 3$, т.е. $\boxed{2, 3}$.}
\end{quote}

Токены рассуждений не видны конечным пользователям, но используются
в процессе обучения: функция верификации извлекает содержимое
окружения \texttt{\textbackslash boxed\{\}} для вычисления вознаграждения.
\end{example}

\subsection{Почему более длинное рассуждение помогает}
\label{subsec:why-cot-helps}

Chain-of-thought работает по информационно-теоретической причине.
Вероятность того, что модель даст правильный финальный ответ $a^*$
при промпте $x$, может быть записана по правилу цепочки вероятностей как
\begin{equation}
\label{eq:cot-chain-rule}
\Prob(a^* \mid x)
  = \sum_{z} \Prob(a^* \mid x, z)\,\Prob(z \mid x).
\end{equation}
Маргинализация по промежуточным следам рассуждений $z$ эквивалентна
их использованию в качестве скрытой переменной. Модель, способная генерировать широкое
разнообразие высококачественных следов рассуждений, эффективно выполняет
неявный поиск по путям решений, прежде чем остановиться на ответе.
Чем длиннее и разнообразнее рассуждение, тем тоньше неявная
сетка поиска.

Эта связь с \emph{поиском}\index{inference-time search} не просто
метафорична. Её можно формализовать, определив качество
следа рассуждений $z$ для промпта $x$ и ответа $a^*$ как
\[
q(z; x, a^*) = \Prob(a^* \mid x, z).
\]
Эффективная точность модели равна $\E_z[q(z; x, a^*)]$, где
математическое ожидание берётся по распределению модели над следами. Более длинный
контекстный буфер допускает больше информации в $z$, что повышает
потолок $q(z; x, a^*)$ для сложных задач. Эмпирически на
олимпиадных математических тестах корреляция между средним количеством
выходных токенов и точностью сильна и монотонна на нескольких
порядках величины вычислений\,---\,масштабный закон на этапе вывода
(Раздел~\ref{sec:inference-scaling}).

\subsection{Токены рассуждений и формат}
\label{subsec:thinking-tokens}

Современные модели рассуждений формализуют chain-of-thought с помощью специальных
\emph{токенов рассуждений}\index{thinking tokens}. Генерация
структурирована следующим образом:
\begin{equation}
\label{eq:thinking-format}
y = \underbrace{\texttt{<thinking>}\; z_1 z_2 \cdots z_T
                \;\texttt{</thinking>}}_{\mbox{след рассуждений, не показывается пользователю}}
    \;\;\underbrace{a_1 a_2 \cdots a_K}_{\mbox{финальный ответ}}.
\end{equation}
Последовательность токенов $z_1, \ldots, z_T$ иногда называется
\emph{черновиком}\index{scratchpad}. Для простого арифметического
вопроса $T$ может составлять несколько десятков токенов; для сложной олимпиадной
задачи $T$ может достигать десятков тысяч токенов, что соответствует
нескольким минутам генерации на современном оборудовании.

\begin{remark}[Длинный контекст как предпосылка]
Обучение моделей рассуждений требует, чтобы стратегия генерировала и обрабатывала
очень длинные последовательности. Практическим катализатором этого сдвига стала
разработка эффективных механизмов внимания (например, FlashAttention)
и контекстных буферов размером от $128\text{k}$ до $1\text{M}$ токенов. Без этих
инженерных достижений описанный в
Разделе~\ref{sec:training-reasoning} режим обучения не был бы вычислительно
осуществим.
\end{remark}

\subsection{Связь с вычислениями на этапе вывода}
\label{subsec:cot-inference-compute}

Расширенный формат генерации~\eqref{eq:thinking-format}
напрямую связывает обучение с распределением вычислений на этапе вывода. Определим
\emph{токенный бюджет}\index{token budget} $B$ как максимальное число
токенов, которые модель может сгенерировать до получения финального ответа. Для
фиксированной стратегии $\pi_\theta$ точность на наборе задач является
неубывающей функцией от $B$:
\[
\text{Acc}(B) = \Prob_{y \sim \pi_\theta(\cdot|x),\,|y| \le B}
               \bigl[\textsc{Verify}(y, a^*) = \textsc{True}\bigr].
\]
Обучение с RLVR на неограниченной генерации учит модель
\emph{самостоятельно распределять} токенный бюджет: сложные задачи порождают более длинные
следы, простые разрешаются быстро. Это адаптивное распределение является
ключевым практическим преимуществом перед альтернативами с фиксированными вычислениями.

% ===========================================================
\section{Обучение моделей рассуждений}
\label{sec:training-reasoning}
% ===========================================================

\subsection{Двухэтапный рецепт}

Доминирующий рецепт обучения моделей рассуждений состоит из двух
этапов, проиллюстрированных здесь на примере DeepSeek R1 как канонического открытого
примера.

\begin{description}[leftmargin=!,labelwidth=3.5cm]
\item[Этап~1: Холодный старт SFT] Базовая предобученная модель сначала
  дообучается с помощью обучения с учителем на относительно небольшом наборе данных
  высококачественных демонстраций chain-of-thought. Этот этап обучает модель
  базовому формату\,---\,использованию тегов \texttt{<thinking>},
  организации многошаговых рассуждений,\,---\,и гарантирует, что
  последующее RL будет иметь хорошо обусловленное начальное распределение для возмущений.
  Без этого холодного старта обучение RLVR от базовой модели может расходиться
  или производить искажённые следы рассуждений.

\item[Этап~2: Крупномасштабный RLVR] Модель, инициализированная через SFT, затем
  обучается с алгоритмом градиента стратегии (обычно GRPO или его
  вариантом; см. Раздел~\ref{sec:practical-reasoning}) на большом
  наборе данных верифицируемых задач. В отличие от RLHF, который обычно выполняется
  $O(1)$ эпох по данным предпочтений, RLVR возвращается к одним и тем же
  обучающим задачам сотни или тысячи раз, позволяя
  стратегии обнаружить и закрепить верные паттерны рассуждений, которые были
  лишь латентными в базовой модели.
\end{description}

\begin{example}[Рецепт обучения DeepSeek R1]
\label{ex:deepseek-r1}
\index{DeepSeek R1}
DeepSeek R1~(январь 2025) обучал базовую модель с приблизительно
600 миллиардами параметров на учебном плане, смешивавшем задачи математического
рассуждения (преимущественно из наборов данных олимпиадной математики),
задачи синтеза кода с верификацией модульными тестами и общие
задачи следования инструкциям. Этап RL использовал GRPO с бинарными
вознаграждениями от функции верификации в сочетании с небольшими вспомогательными
наградами за формат для обеспечения последовательного использования структуры \texttt{<thinking>}.
Впоследствии модель была дистиллирована до меньших размеров
(7B--70B) путём дообучения на выходах полной модели, что сделало
сильные возможности рассуждений доступными на потребительском оборудовании.
\end{example}

\subsection{Фильтрация по сложности}
\label{subsec:difficulty-filtering}

Тонким, но критически важным аспектом обучения RLVR является выбор того, какие
задачи включить в обучающий набор. Вспомним обновление градиента стратегии:
градиент $J_{\mathrm{RLVR}}$ по $\theta$
ненулевой только тогда, когда разные развёртки для одного промпта получают
разные вознаграждения. Если текущая стратегия решает задачу в 100\% случаев,
все развёртки получают вознаграждение~$1$, преимущество тождественно
равно нулю, и градиент исчезает. Аналогично, если стратегия проваливается в 100\%
случаев, все развёртки получают вознаграждение~$0$ и не генерируется никакого градиента.\index{difficulty filtering}

\begin{definition}[Коэффициент решаемости задачи]
\label{def:solvability}
Для задачи $x$ с эталоном $a^*$ и стратегией $\pi_\theta$
\textbf{коэффициент решаемости} равен
\[
\sigma(x) = \Prob_{y \sim \pi_\theta(\cdot|x)}
            \bigl[\textsc{Verify}(y, a^*) = \textsc{True}\bigr],
\]
оцениваемый путём сэмплирования $N$ независимых развёрток (обычно $N = 16$).
\end{definition}

Эмпирически задачи с $\sigma(x) \in [0.2, 0.8]$ дают
наиболее сильные градиенты. Задачи вне этого диапазона\,---\,слишком лёгкие или
слишком сложные\,---\,исключаются из обучающего набора до начала RLVR.
Эта \emph{офлайн-фильтрация по сложности} сейчас является стандартным
шагом в конвейерах моделей рассуждений, используемым в DeepSeek R1, Phi-4
Reasoning, INTELLECT-2, Skywork OR-1 и других.

\subsection{Удаление штрафа KL}
\label{subsec:no-kl}

Стандартный RLHF использует штраф KL из~\eqref{eq:rlhf-objective-ch15}, чтобы
предотвратить схлопывание стратегии к выходам с взломом вознаграждения. В
RLVR бинарное вознаграждение нельзя взломать, повторяя фиксированный выход,
поскольку функция верификации детерминирована и специфична для задачи.
В результате многие практики RLVR уменьшают $\beta$ до нуля, полностью удаляя
ограничение KL. Это позволяет стратегии исследовать более свободно
и обнаруживать паттерны рассуждений, отсутствовавшие в
референсной модели. Модели, удалившие штраф KL, включают
Magistral (Mistral, 2025), OpenReasoner-Zero и Skywork OR-1.

\begin{remark}[Удаление KL и сдвиг распределения]
Удаление штрафа KL действительно вводит риск: стратегия может сместиться
к распределению, при котором сам верификатор ненадёжен (например,
если модель научится извлекать ответы из позиций, которые регулярное
выражение не соответствует). Тщательная разработка промптов и тестирование
устойчивости верификатора поэтому важнее в режимах без KL, чем в
стандартном RLHF.
\end{remark}

\subsection{Награды за формат и штрафы за длину}
\label{subsec:format-rewards}

В дополнение к бинарному вознаграждению за задачу практики обычно добавляют
небольшие вспомогательные вознаграждения:

\begin{itemize}[leftmargin=*]
\item \textbf{Награды за формат}: Небольшое положительное вознаграждение ($\approx 0.1$)
  даётся, когда ответ корректно использует структуру \texttt{<thinking>}
  \ldots\ \texttt{</thinking>} перед финальным ответом. Это предотвращает
  смещение модели к формату, который трудно разобрать.

\item \textbf{Штрафы за длину}: Следы рассуждений, более длинные, чем
  необходимо, потребляют лишние вычисления на этапе вывода без улучшения точности.
  Штраф, пропорциональный $\max(0, |z| - L_{\max})$, препятствует
  патологическому обдумыванию\index{overthinking} при сохранении
  возможности модели использовать длинные следы когда это действительно необходимо. Некоторые
  рецепты (например, Kimi 1.5) используют прогрессивный учебный план по длине:
  модель сначала учится рассуждать в жёстком токенном бюджете, а
  бюджет расширяется по мере продвижения обучения.

\item \textbf{Награды за языковую согласованность}: Многоязычные модели могут
  перейти к смешению языков в середине следа (например, начав на английском
  и завершив на китайском), что ухудшает пользовательский опыт. Бинарная
  награда за языковую согласованность используется в DeepSeek R1 и Magistral.
\end{itemize}

\subsection{Нормализация функции потерь и смещение GRPO}
\label{subsec:loss-norm}

Алгоритм GRPO (Глава~\ref{ch:grpo}) вычисляет преимущества путём
нормализации вознаграждений внутри каждой группы развёрток для одного
промпта. Тонким следствием является то, что вклад задачи в
общую функцию потерь не зависит от числа развёрток, сгенерированных
для неё. Если более сложные задачи систематически требуют больше токенов,
нормализация на уровне группы может придавать избыточный вес лёгким коротким задачам
и недостаточный вес сложным. Несколько рецептов моделей
рассуждений (включая Magistral и MiMo) решают это путём нормализации преимуществ
на уровне \emph{батча}, а не группы, или взвешивания функции потерь
по количеству токенов.

\subsection{Самообучаемое рассуждение и дистилляция}
\label{subsec:star-distill}

Два важных варианта стандартного рецепта отходят от чистого RL.

\textbf{STaR}\index{STaR} (Self-Taught Reasoner; Zeiler et al., 2022)
и его потомок \textbf{Quiet-STaR}\index{Quiet-STaR} аппроксимируют
градиент стратегии путём чередования: (1) сэмплирование следов рассуждений
из текущей модели, (2) фильтрация до тех, которые приводят к правильным
ответам, и (3) дообучение модели на отфильтрованных следах с
функцией потерь кросс-энтропии. Это эквивалентно упрощённой форме RL, где
вознаграждение бинарно, а базовая линия равна нулю. STaR аппроксимирует
градиент~\eqref{eq:rlvr-objective} без явного вычисления преимущества, ценой
некоторой статистической эффективности.
Quiet-STaR добавляет инновацию генерации токенов рассуждений перед
каждым предсказанием токена, а не только в начале ответа,
что обеспечивает более богатый обучающий сигнал.

\textbf{Дистилляция}\index{distillation!reasoning} переносит
способность к рассуждению от большой, обученной с RLVR модели к меньшей
модели-ученику посредством обучения с учителем на выходах
chain-of-thought модели-учителя. Поскольку ученик учится на завершённых следах
рассуждений учителя\,---\,включая токены рассуждений,\,---\,он может
достичь производительности рассуждений намного выше той, которую прямое RLVR
на самом ученике позволило бы достичь. DeepSeek R1 опубликовал
дистиллированные контрольные точки при 7B и 70B, которые превзошли свои аналоги
того же размера без дистилляции на олимпиадной математике.

% ===========================================================
\section{Масштабирование на этапе вывода}
\label{sec:inference-scaling}
% ===========================================================

\subsection{Гипотеза о вычислениях во время тестирования}

Масштабирование на этапе вывода\index{inference-time scaling}~--- это
эмпирическое наблюдение, что для фиксированной обученной стратегии выделение большего
объёма вычислений во время тестирования\,---\,путём генерации большего числа токенов, большего числа
ответов или и того, и другого,\,---\,приводит к измеримым улучшениям точности задач.
Это отличается от масштабирования на этапе обучения (более крупные модели
или больше обучающих данных): это масштабирование при развёртывании.

Пусть $\mathcal{P}$ обозначает набор тестовых задач, а
$C$~--- бюджет вычислений, измеренный в FLOPs (или, эквивалентно, в
суммарном числе сгенерированных выходных токенов). Закон масштабирования вывода утверждает, что
существует функция $f$ такая, что
\begin{equation}
\label{eq:inference-scaling}
\text{Acc}(C) \approx f(C) = f_0 + \alpha \log C + O((\log C)^{-1}),
\end{equation}
для констант $f_0$ и $\alpha > 0$. Эмпирически это логарифмически-линейное
соотношение справедливо на двух-трёх порядках величины вычислений во время тестирования
на олимпиадных математических тестах, прежде чем насыщается по мере исчерпания задач.

\begin{figure}[ht]
\centering
\begin{tikzpicture}[xscale=1.5, yscale=0.07]
  % Сетка
  \draw[gray!30, dashed] (0,30) grid[xstep=1,ystep=10] (6,90);
  % Оси
  \draw[->] (0,30) -- (6.3,30);
  \draw[->] (0,30) -- (0,95) node[above, font=\small]{Точность (\%)};
  % Подпись оси X отдельно, чтобы не вылезала за страницу
  \node[font=\small, align=center] at (3.2,20)
    {Вычисления во время тестирования};
  % Метки Y
  \foreach \y in {30,40,50,60,70,80,90}
    \draw (-0.1,\y) node[left, font=\scriptsize]{\y} -- (0.05,\y);
  % Метки X
  \foreach \x/\lab in {1/$10^1$,2/$10^2$,3/$10^3$,4/$10^4$,5/$10^5$}
    \draw (\x,29) node[below, font=\scriptsize]{\lab} -- (\x,30.5);
  % Единичная развёртка
  \draw[thick, blue] plot[smooth, domain=0.2:5.8, samples=40]
    (\x, {32 + 38*(1 - exp(-0.7*\x))});
  % Голосование большинства
  \draw[thick, red, dashed] plot[smooth, domain=0.2:5.8, samples=40]
    (\x, {50 + 40*(1 - exp(-0.55*\x))});
  % Сэмплирование с направлением по ценности
  \draw[thick, orange, densely dotted] plot[smooth, domain=0.2:5.8, samples=40]
    (\x, {58 + 35*(1 - exp(-0.5*\x))});
  % Линия насыщения
  \draw[gray, dashed] (5.2,30) -- (5.2,92)
    node[above, font=\tiny, gray]{насыщение};
  % Легенда
  \draw[thick, blue] (3.5,42) -- ++(0.5,0)
    node[right, font=\scriptsize]{Единичная развёртка (длинный CoT)};
  \draw[thick, red, dashed] (3.5,38) -- ++(0.5,0)
    node[right, font=\scriptsize]{Голосование большинства ($k$ развёрток)};
  \draw[thick, orange, densely dotted] (3.5,34) -- ++(0.5,0)
    node[right, font=\scriptsize]{Сэмплирование с направлением по ценности};
\end{tikzpicture}
\caption{Схематичные кривые масштабирования на этапе вывода для олимпиадной
  математики. Точность растёт логарифмически-линейно с вычислениями во время тестирования для
  ряда стратегий. Голосование большинства по $k$ параллельным развёрткам
  даёт более высокую точность при тех же вычислениях, чем единственная длинная
  развёртка; сэмплирование с направлением по ценности извлекает дополнительные улучшения, используя
  процессную модель вознаграждения для управления поиском по лучу. Все кривые насыщаются,
  как только достигается потолок возможностей модели.}
\label{fig:inference-scaling}
\end{figure}

\subsection{Голосование большинства}
\label{subsec:majority-voting}

Простейшей стратегией масштабирования на этапе вывода является \emph{голосование
большинства}\index{majority voting} (также называемое \emph{самосогласованностью}):
сэмплируем $k$ независимых ответов $y_1, \ldots, y_k$ из стратегии,
извлекаем финальный ответ $a_i$ из каждого и возвращаем ответ с наибольшим числом голосов:
\begin{equation}
\label{eq:majority-vote}
\hat{a} = \argmax_{a} \sum_{i=1}^{k} \ind\{a_i = a\}.
\end{equation}
Если отдельные ответы верны с вероятностью $p >
\frac{1}{2}$, голосование большинства даёт верный ответ с вероятностью, которая
стремится к~$1$ экспоненциально быстро при росте $k$ (следствие
закона больших чисел). На практике $k = 8$--$64$ развёрток
достаточно для большинства олимпиадных тестов.

\begin{proposition}[Граница точности голосования большинства]
\label{prop:majority-vote}
Предположим, что каждая из $k$ независимых развёрток верна с вероятностью
$p \in (\frac{1}{2}, 1)$ и существует не более $m$ возможных ответов.
Голосование большинства ошибается с вероятностью не более
\[
\Prob(\hat{a} \ne a^*) \le
\binom{k}{\lfloor k/2 \rfloor} p^{\lfloor k/2 \rfloor}(1-p)^{k - \lfloor k/2 \rfloor}
\;\le\; \exp\bigl(-k\,D_{\mathrm{KL}}(1/2 \,\|\, 1-p)\bigr),
\]
где $D_{\mathrm{KL}}$ обозначает бинарную KL-дивергенцию. Ошибка
убывает экспоненциально при росте $k$.
\end{proposition}

\subsection{Сэмплирование с направлением по ценности}
\label{subsec:value-guided}

Голосование большинства считает все развёртки одинаково правдоподобными. Более
сложная стратегия использует \emph{процессную модель вознаграждения}
(PRM)\index{process reward model}, которая присваивает оценки
\emph{промежуточным} шагам рассуждения, а не только финальному ответу.
Дан частичный след $z_{1:t}$, PRM предсказывает вероятность
того, что рассуждение приведёт к верному финальному ответу. Эта оценка
может направлять генерацию посредством поиска по лучу или выбора лучшего из $k$:
\begin{equation}
\label{eq:prm-selection}
\hat{y} = \argmax_{y \in \{y_1,\ldots,y_k\}} \mathrm{PRM}(y).
\end{equation}
Сэмплирование с направлением по ценности последовательно превосходит голосование большинства при равных
вычислениях, поскольку концентрирует бюджет сэмплирования на перспективных путях
рассуждений, а не усредняет по всем им. Накладные расходы~--- стоимость
запуска скорера PRM, который обычно значительно меньше модели
стратегии.

\subsection{Вычислительно-оптимальный вывод}
\label{subsec:compute-optimal}

Естественный вопрос~--- как разделить фиксированный бюджет вычислений $C$ между
длиной отдельных следов (больше токенов рассуждений на развёртку)
и числом развёрток $k$. Определяя стоимость одной развёртки как
$c(B) \propto B$ (где $B$~--- токенный бюджет на развёртку), общие
вычисления составляют $C = k \cdot c(B)$. Вычислительно-оптимальная стратегия
решает:
\[
(k^*, B^*) = \argmax_{k,\,B:\,k\cdot B = C}\; \text{Acc}(k, B),
\]
где $\text{Acc}(k, B)$~--- точность при голосовании большинства с $k$
развёртками длины $B$. Эмпирически этот компромисс отдаёт предпочтение более длинным
отдельным следам для сложных задач (где $B^*$ велик) и большему числу
коротких развёрток для лёгких задач (где $k^*$ велик).

% ===========================================================
\section{Использование инструментов и вызов функций}
\label{sec:tool-use}
% ===========================================================

\subsection{Мотивация}

Языковые модели, обученные исключительно на тексте, имеют фундаментальное ограничение:
они не могут получать доступ к информации за пределами своих обучающих данных, выполнять
код или взаимодействовать с внешними сервисами. Использование инструментов\index{tool use}
расширяет пространство действий модели от генерации токенов до
\emph{вызовов функций}: структурированных выходов, запрашивающих выполнение
внешнего инструмента и встраивающих результат в последующую генерацию.
Это качественное расширение парадигмы генерации:
\begin{equation}
\label{eq:tool-loop}
x \;\longrightarrow\; (z_1, \underbrace{f_1(q_1)}_{\mbox{вызов инструмента}},
r_1, z_2, \underbrace{f_2(q_2)}_{\mbox{вызов инструмента}}, r_2, \ldots, a),
\end{equation}
где $f_i$~--- инструмент (например, API веб-поиска или интерпретатор Python),
$q_i$~--- запрос, который модель отправляет ему, $r_i$~--- возвращённый результат,
а генерация модели чередует рассуждение ($z_i$)
с вызовами инструментов и их результатами.

\subsection{Формат вызова функций}
\label{subsec:function-calling-format}

На практике вызовы инструментов выражаются как структурированные объекты JSON,
встроенные в поток генерации. Типичный формат:

\begin{example}[Вызов инструмента в формате JSON]
\label{ex:tool-call-json}
\index{function calling!format}
Рассмотрим модель, отвечающую на вопрос: \emph{«Какова нынешняя численность населения
Франции?»} Генерация может выглядеть следующим образом:

\begin{verbatim}
<thinking>
Мне нужно найти актуальную численность населения Франции.
</thinking>
{
  "tool": "web_search",
  "query": "численность населения Франции 2025"
}
[Результат инструмента: «Население Франции: около 68,4 миллиона (2025)»]
Нынешняя численность населения Франции составляет около 68,4 миллиона человек.
\end{verbatim}

Объект JSON анализируется средой выполнения, инструмент
вызывается, а результат добавляется в контекст модели до
продолжения генерации. Модель никогда не видит внутренние механизмы сырого API; она
взаимодействует с чистым интерфейсом именованных инструментов с типизированными параметрами.
\end{example}

Модель должна одновременно освоить три возможности: (1) распознавать,
когда нужен вызов инструмента, (2) генерировать корректный JSON,
соответствующий схеме инструмента, и (3) правильно включать
результат инструмента в последующее рассуждение.

\subsection{Обучение использованию инструментов}
\label{subsec:training-tool-use}

Использование инструментов обычно обучается в два этапа:

\begin{description}[leftmargin=!,labelwidth=3.5cm]
\item[SFT на демонстрациях] Модель сначала дообучается на
  кураторском наборе данных траекторий использования инструментов: тройках (промпт, последовательность
  вызовов инструментов, финальный ответ), собранных из демонстраций людей
  или больших моделей. Это обучает модель формату и
  базовому решению о том, когда вызывать инструмент.

\item[RL для многошагового использования инструментов] Одного SFT недостаточно для
  сложных многошаговых сценариев использования инструментов, поскольку правильная последовательность
  вызовов инструментов зависит от результатов предыдущих вызовов способом, который
  трудно перечислить в фиксированном наборе данных демонстраций. RL позволяет
  модели обнаруживать эффективные стратегии вызова инструментов путём
  взаимодействия с живой средой инструментов и получения бинарного
  вознаграждения за выполнение задачи.
\end{description}

Сигнал вознаграждения для RL с использованием инструментов обычно основан на результате: модель
получает вознаграждение~$1$, если финальный ответ (после всех вызовов инструментов)
верен, и~$0$ иначе. Это в точности фреймворк RLVR
из Раздела~\ref{sec:rlvr-ch15}, теперь расширенный на многошаговые траектории.
Средой является сама экосистема инструментов.

\subsection{SWiRL: пошаговый RL для использования инструментов}
\label{subsec:swirl}

\index{SWiRL}
\textbf{SWiRL} (Step-Wise Reinforcement Learning~--- пошаговое обучение с подкреплением)~--- это фреймворк,
разработанный специально для многошаговых задач использования инструментов, где вознаграждения за результат
разрежены. Вместо того чтобы ждать конца траектории для присвоения вознаграждения,
SWiRL предоставляет \emph{процессные вознаграждения} на каждом вызове инструмента.
Пошаговое вознаграждение $r_t$ вычисляется процессной моделью вознаграждения, которая
оценивает, был ли $t$-й вызов инструмента полезным и корректно сформированным.
Общее вознаграждение:
\begin{equation}
\label{eq:swirl-reward}
R = \sum_{t=1}^{T} \gamma^{t-1} r_t + \lambda \cdot r_{\mathrm{outcome}},
\end{equation}
где $r_{\mathrm{outcome}} \in \{0, 1\}$~--- бинарное вознаграждение за выполнение задачи,
а $\lambda$~--- гиперпараметр взвешивания. Пошаговые вознаграждения
снижают дисперсию оценки градиента стратегии, обеспечивая более
частую обратную связь, ценой необходимости в пошаговом оценщике.

% ===========================================================
\section{Многошаговые агенты и агентное RL}
\label{sec:agentic-rl}
% ===========================================================

\subsection{Цикл агента}

Когда использование инструментов охватывает множество шагов и модель должна планировать через
длинный горизонт, мы вступаем в режим \emph{агентного RL}\index{agentic RL}.
Модель теперь функционирует как полноценный агент RL, взаимодействующий со
средой $\mathcal{E}$ (набором доступных инструментов и
внешних систем). Стандартный цикл агента:

\begin{enumerate}
\item \textbf{Наблюдение}: Получить наблюдение $o_t$ из
  среды (например, содержимое веб-страницы, вывод
  терминальной команды или ответ API).
\item \textbf{Рассуждение}: Сгенерировать след рассуждений $z_t$, обусловленный
  на полном контексте $(x, o_1, a_1, \ldots, o_{t-1}, a_{t-1}, o_t)$.
\item \textbf{Действие}: Вывести действие $a_t$, которым может быть вызов инструмента,
  навигационная команда или текстовый вывод.
\item \textbf{Наблюдение}: Получить $o_{t+1}$ из среды. Перейти к~2.
\end{enumerate}

Это в точности цикл агента MDP из Главы~\ref{ch:mdp}, теперь
реализованный с языковой моделью в качестве стратегии и программной
средой в качестве функции перехода.

\begin{figure}[ht]
\centering
% Упрощённый цикл агентного RL: понятная разметка с метками вдали от стрелок
\begin{tikzpicture}[
  node distance=1.4cm and 2.6cm,
  box/.style={rectangle, rounded corners=4pt, draw, thick,
              minimum width=2.4cm, minimum height=0.8cm,
              font=\small, align=center},
  grp/.style={rectangle, rounded corners=6pt, draw, thick,
              minimum width=2.8cm, minimum height=2.8cm,
              font=\small, align=center},
  arrow/.style={-{Stealth[length=6pt]}, thick},
  darrow/.style={-{Stealth[length=6pt]}, thick, dashed},
  lbl/.style={font=\scriptsize, fill=white, inner sep=1pt},
]

% ---- Группа агента (левый столбец) ----
\node[grp, fill=blue!8] (agentbox) {};
\node[font=\small\bfseries, above=0.05cm of agentbox.north]
      {LLM-агент $\pi_\theta$};
\node[box, fill=blue!18, minimum width=2.0cm, minimum height=0.6cm,
      above=0.05cm of agentbox.center] (think) {Рассуждение $z_t$};
\node[box, fill=blue!18, minimum width=2.0cm, minimum height=0.6cm,
      below=0.05cm of agentbox.center] (act) {Действие $a_t$};

% ---- Группа среды (правый столбец) ----
\node[grp, fill=orange!8, right=2.6cm of agentbox] (envbox) {};
\node[font=\small\bfseries, above=0.05cm of envbox.north]
      {Среда};
\node[box, fill=orange!18, minimum width=2.0cm, minimum height=0.6cm,
      above=0.05cm of envbox.center] (tools) {Инструменты и API};
\node[box, fill=orange!18, minimum width=2.0cm, minimum height=0.6cm,
      below=0.05cm of envbox.center] (envobs) {Набл.\  $o_{t+1}$};

% ---- Задача (над агентом) ----
\node[box, fill=green!10, above=1.2cm of agentbox, minimum width=2.6cm]
      (task) {Задача / Промпт $x$};

% ---- Контекстный буфер (под агентом) ----
\node[box, fill=gray!12, below=1.2cm of agentbox, minimum width=2.6cm]
      (ctx) {Контекстный буфер};

% ---- Вознаграждение (под средой) ----
\node[box, fill=red!10, below=1.2cm of envbox, minimum width=2.6cm]
      (reward) {Вознаграждение $r$};

% ---- Стрелки ----
% Задача -> Агент
\draw[arrow] (task.south) -- (agentbox.north);

% Действие -> Инструменты
\draw[arrow] ([yshift=-2pt]act.east) -- ([yshift=-2pt]tools.west)
      node[lbl, below, midway] {действие $a_t$};

% Набл -> Рассуждение
\draw[arrow] ([yshift=2pt]envobs.west) -- ([yshift=2pt]think.east)
      node[lbl, above, midway] {набл.\  $o_{t+1}$};

% Агент -> Контекст
\draw[arrow] (agentbox.south) -- (ctx.north);

% Контекст -> Вознаграждение
\draw[arrow] (ctx.east) -- (reward.west)
      node[lbl, below, midway] {конец эпизода};

% Вознаграждение -> Среда (пунктирная, верификация)
\draw[darrow] (reward.north) -- (envbox.south)
      node[lbl, right, midway] {верификация};

\end{tikzpicture}
\caption{Цикл агентного RL. LLM-агент $\pi_\theta$ наблюдает
  среду, генерирует след рассуждений и производит действие (вызов
  инструмента или текстовый вывод). Среда выполняет действие и
  возвращает наблюдение. Полный контекст накапливается в контекстном
  буфере модели. После завершения эпизода вычисляется сигнал вознаграждения
  (на основе результата или пошагового) и используется для обновления градиента стратегии.}
\label{fig:agent-loop}
\end{figure}

\subsection{Задачи и среды агентного RL}
\label{subsec:agentic-tasks}

Агентное RL было применено к растущему числу задач:

\begin{description}[leftmargin=!,labelwidth=3.5cm]
\item[Браузинг веба] Агент навигирует по веб-браузеру (через Playwright
  или аналогичное), читает страницы, переходит по ссылкам и заполняет формы для
  выполнения задач поиска информации или транзакций. Вознаграждение
  бинарно: цель задачи достигнута или нет.

\item[Выполнение кода] Агент пишет код, выполняет его в
  изолированном интерпретаторе Python, читает вывод и итерирует. Сигнал
  вознаграждения через модульные тесты из Раздела~\ref{sec:rlvr-ch15} применяется
  напрямую.

\item[Взаимодействие с API] Агент вызывает REST API, читает ответы
  и цепочкой API-вызовов завершает многошаговые рабочие процессы (например, бронирование
  события в календаре после проверки доступности). Вознаграждение поступает от
  верификации финального состояния системы.

\item[Кодирование на уровне репозитория] Агент читает и пишет файлы в
  репозитории кода, запускает тесты и устраняет проблемы из отчётов об ошибках.
  SWE-bench\index{SWE-bench}~--- стандартный тест, измеряющий
  долю разрешённых агентом задач с GitHub.
\end{description}

\subsection{Отнесение заслуг в многошаговых траекториях}
\label{subsec:credit-agentic}

Фундаментальный вызов агентного RL\index{credit assignment!agentic}~---
отнесение заслуг: когда траектория из $T$ шагов получает единственное
конечное вознаграждение, каждое из $T$ действий должно быть приписано соответствующим образом.
Это проблема отложенного вознаграждения из классического RL
(Глава~\ref{ch:td}), усиленная длиной траекторий языковых моделей.

Два подхода решают её:

\begin{itemize}[leftmargin=*]
\item \textbf{Вознаграждения за результат}: Начислить каждому токену в траектории
  одинаковое конечное вознаграждение (возможно, дисконтированное). Это просто,
  но имеет высокую дисперсию: ранние действия, настраивающие успех, вознаграждаются
  наравне с единственным финальным токеном, подтверждающим успех.

\item \textbf{Процессные вознаграждения}: Обучить PRM на уровне шагов для присвоения
  вознаграждений $r_t$ на каждом шаге. Это значительно снижает дисперсию,
  но требует дополнительного надзора для обучения PRM. PRM можно
  обучать на разметке людьми того, какие промежуточные шаги
  верны, или (более масштабируемо) с использованием развёрток Монте-Карло из каждого
  промежуточного состояния для оценки вероятности конечного успеха.
\end{itemize}

Компромисс между процессными и результирующими вознаграждениями (рассмотренный подробнее
в Главе~\ref{ch:reward-models}) особенно остёр в
агентной обстановке, где траектории могут охватывать десятки вызовов инструментов
и сотни тысяч токенов.

\subsection{Исследование в агентном RL}
\label{subsec:agentic-exploration}

Классические стратегии исследования в RL (epsilon-жадный, сэмплирование Томпсона,
внутренняя мотивация; Глава~\ref{ch:intro}) приобретают новый
характер, когда стратегией является языковая модель. Априорное знание модели от
предобучения уже кодирует богатое распределение над последовательностями действий,
что обеспечивает неявное исследование. Однако это априорное знание может быть
сосредоточено на безопасных, привычных последовательностях действий, которые избегают
долгосрочных стратегий, необходимых для сложных агентных задач.

Практические подходы к агентному исследованию включают:

\begin{itemize}[leftmargin=*]
\item \textbf{Масштабирование температуры}: Увеличение температуры генерации
  во время обучения разнообразит развёртки, аналогично увеличению
  epsilon в epsilon-жадном методе.

\item \textbf{Учебный план по сложности задач}: Начало с коротких,
  лёгких задач и постепенное введение более длинных, сложных гарантирует,
  что модель получает адекватные сигналы обратной связи на протяжении всего обучения.

\item \textbf{Рандомизация среды}: Изменение экземпляров задач
  (разные веб-страницы, разные кодовые репозитории, разные конечные точки API)
  предотвращает заучивание стратегией конкретных траекторий.
\end{itemize}

% ===========================================================
\section{Практические соображения}
\label{sec:practical-reasoning}
% ===========================================================

\subsection{Выбор алгоритма}
\label{subsec:algorithm-choices}

Таблица~\ref{tab:reasoning-algorithms} обобщает основные алгоритмы RL,
используемые для обучения рассуждению и агентного обучения.

\begin{table}[ht]
\centering
\footnotesize
\caption{Алгоритмы оптимизации стратегии для обучения моделей рассуждений.}
\label{tab:reasoning-algorithms}
\begin{tabular}{@{}lp{3.5cm}p{2.5cm}p{2.5cm}@{}}
\toprule
\textbf{Алгоритм} & \textbf{Механизм} & \textbf{Ключевое свойство}
  & \textbf{Используется в} \\
\midrule
PPO & Отсечённое приближение + сеть ценности
  & Устойчив, хорошо изучен & Tülu~3, Phi-4 \\
GRPO & Групповая норм.\  преимуществ, без сети ценности
  & Эффективен по памяти & DeepSeek R1, Qwen~3 \\
DAPO & Асимметричное отсечение
  & Лучшее исследование & RAGEN, INTELLECT-2 \\
CISPO & Отсечённое взвешивание важности
  & Офлайн-устойчивость & MiniMax-M1 \\
MuonClip & Отсечение QK-нормы
  & Предотвращает скачки функции потерь & Kimi~K2 \\
\bottomrule
\end{tabular}
\end{table}

\textbf{PPO}\index{PPO!reasoning} (Proximal Policy Optimisation~--- проксимальная оптимизация стратегии)
остаётся наиболее широко изученным алгоритмом дообучения RL благодаря
его теоретическим гарантиям и обширному эмпирическому опыту.
Он требует поддержки отдельной сети ценности, которая удваивает
потребление памяти, но снижает дисперсию градиента через обученную базовую линию.

\textbf{GRPO}\index{GRPO} (Group Relative Policy Optimisation~--- групповая относительная оптимизация стратегии) избегает
сети ценности путём вычисления оценок преимущества на уровне промпта из
эмпирического среднего и дисперсии вознаграждений внутри группы из $k$
развёрток:
\begin{equation}
\label{eq:grpo-advantage-ch15}
\hat{A}_i = \frac{r_i - \bar{r}}{\sigma_r}, \qquad
\bar{r} = \frac{1}{k}\sum_{j=1}^k r_j, \quad
\sigma_r = \sqrt{\frac{1}{k}\sum_{j=1}^k (r_j - \bar{r})^2}.
\end{equation}
При бинарных вознаграждениях RLVR, $r_i \in \{0, 1\}$, это упрощается до
$\hat{A}_i = +\sigma_r^{-1}$ для правильных завершений и
$\hat{A}_i = -\bar{r}/\sigma_r$ для неправильных. GRPO стал
доминирующим в обучении моделей рассуждений, поскольку эффективен по памяти
и прост в реализации.

\textbf{DAPO}\index{DAPO} (Decoupled Clip and Dynamic Sampling Policy
Optimisation~--- оптимизация стратегии с разделённым отсечением и динамическим сэмплированием) модифицирует GRPO, используя асимметричные границы отсечения: большее
отсечение на положительной стороне позволяет более агрессивные обновления в сторону
высоковознаграждаемых ответов, тогда как отсечение на отрицательной стороне остаётся консервативным.
Это способствует исследованию правильных стратегий рассуждения, которые
текущая стратегия достигает лишь редко.

\subsection{Онлайн- и офлайн-RL}
\label{subsec:online-offline}

\begin{description}[leftmargin=!,labelwidth=3.5cm]
\item[Онлайн RL] Развёртки генерируются из текущей стратегии
  $\pi_\theta$ и немедленно используются для обновления градиента. Это поддерживает
  распределение обучения on-policy, что теоретически чисто,
  но практически медленно: стратегия должна ждать завершения генерации
  перед обновлением.

\item[Офлайн RL] Развёртки предварительно генерируются (из более ранней
  контрольной точки стратегии или из другой модели) и хранятся в
  буфере воспроизведения. Обновления используют веса важности для коррекции
  несоответствия распределений. Офлайн RL значительно улучшает
  пропускную способность, но вводит офлайн-смещение, если буфер воспроизведения устарел.

\item[Асинхронный RL] Практический компромисс, используемый несколькими
  рецептами моделей рассуждений: генерация выполняется асинхронно на одном наборе
  GPU, пока обновления градиента выполняются на другом. Это поддерживает
  стратегию близкой к on-policy, избегая простоя вычислений.
\end{description}

\subsection{Инфраструктура в масштабе}
\label{subsec:infra}

Обучение модели рассуждений требует инфраструктуры, отличающейся от
стандартного предобучения языковых моделей в нескольких аспектах:

\begin{itemize}[leftmargin=*]
\item \textbf{Пропускная способность генерации}: Обучение RLVR требует сэмплирования
  $k = 8$--$16$ развёрток на промпт, каждая потенциально тысячи
  токенов длиной. Необходима эффективная батчевая генерация с непрерывным батчингом
  (как в vLLM).

\item \textbf{Верификация в масштабе}: Верификатор (регулярный экстрактор,
  изолированная среда для кода) должен работать параллельно для тысяч завершений в
  секунду, не становясь узким местом. Верификация кода
  особенно дорогостояща, так как требует изоляции процессов.

\item \textbf{Управление памятью}: Устранение GRPO сети ценности
  вдвое уменьшает потребление памяти GPU по сравнению с PPO, но хранение
  $k$ полных последовательностей развёрток на промпт всё равно требует тщательного планирования памяти.

\item \textbf{Нормализация вознаграждений и журналирование}: Поскольку вознаграждения RLVR
  бинарны и разрежены в начале обучения, необходим тщательный мониторинг
  скользящего среднего вознаграждения, длины развёртки и соответствия формату
  для обнаружения нестабильностей обучения до их каскадного развития.
\end{itemize}

\subsection{Инструменты с открытым исходным кодом}
\label{subsec:open-source-tooling}

Несколько фреймворков с открытым исходным кодом сделали обучение RLVR доступным
за пределами хорошо финансируемых промышленных лабораторий:

\begin{description}[leftmargin=!,labelwidth=2.5cm]
\item[TRL] Библиотека Transformer Reinforcement Learning от Hugging Face
  реализует PPO, GRPO и DPO с первоклассной поддержкой
  верифицируемых вознаграждений. Её \texttt{GRPOTrainer} интегрируется напрямую
  с экосистемами \texttt{datasets} и \texttt{transformers}.

\item[veRL] Высокопропускной фреймворк обучения RL от Volcengine,
  поддерживающий асинхронные конвейеры генерации-обновления и
  достигающий максимальной пропускной способности обучения для крупных моделей.

\item[OpenRLHF] Распределённый фреймворк RLHF и RLVR, построенный на
  DeepSpeed и Ray, поддерживающий модели до сотен миллиардов
  параметров с гибкими API функций вознаграждения.
\end{description}

% ===========================================================
\section{Итоги и открытые вопросы}
\label{sec:reasoning-summary}
% ===========================================================

В этой главе прослежена дуга от RLHF\,---\,где обученная
модель вознаграждения обеспечивала несовершенные сигналы предпочтений,\,---\,к RLVR,
где детерминированные функции верификации обеспечивают агрессивное,
долгосрочное обучение RL на верифицируемых областях. Ключевые концептуальные
вклады:

\begin{itemize}[leftmargin=*]
\item \textbf{Верифицируемые вознаграждения} устраняют модель вознаграждения как
  узкое место и источник отказов. Бинарные вознаграждения от проверочных функций
  не содержат шума и не могут быть сверхоптимизированы.

\item \textbf{Рассуждения типа chain-of-thought} связывают генерацию токенов
  с неявным поиском: более длинные следы рассуждений реализуют более тонкий
  поиск по путям решений.

\item \textbf{Масштабирование на этапе вывода} показывает, что вычисления можно
  продуктивно использовать при развёртывании, не только при обучении. Голосование
  большинства и сэмплирование с направлением по ценности извлекают улучшения точности из
  дополнительной генерации без повторного обучения.

\item \textbf{Использование инструментов и агентное RL} расширяют фреймворк на
  интерактивные среды, превращая языковые модели из пассивных
  генераторов текста в активных агентов решения проблем.
\end{itemize}

\paragraph{Открытые вопросы.} Область развивается быстро, и ряд
фундаментальных вопросов остаётся нерешённым:

\begin{itemize}[leftmargin=*]
\item \textbf{Обобщение RLVR}: Переносятся ли навыки рассуждения, освоенные
  на математике и коде, на неверифицируемые области
  (открытое письмо, стратегическое рассуждение, безопасность)? Ранние данные
  указывают на частичный перенос, но механизм неясен.

\item \textbf{Пределы масштабирования}: Как далеко распространяются законы масштабирования
  на этапе вывода? В какой момент генерация большего числа токенов даёт убывающую
  отдачу независимо от мощности модели?

\item \textbf{Процессные модели вознаграждения в масштабе}: Обучение надёжной
  PRM требует надзора на уровне шагов, дорогостоящего для сбора.
  Оценка Монте-Карло корректности шагов шумна. Лучшие методы
  обучения PRM~--- активная область исследований.

\item \textbf{Безопасность агентных систем}: Агент, способный просматривать
  веб, выполнять код и вызывать API, имеет значительно большее пространство действий,
  чем текстовый чат-бот. Согласование целей
  многошаговых агентов с человеческими ценностями~--- открытый вызов, который
  классический RLHF не был разработан для решения.

\item \textbf{Взаимодействие с предобучением}: Ряд результатов указывает на то,
  что определённый уровень возможностей предобучения является предпосылкой
  для работы RLVR. Понимание порога и того, как его снизить,
  демократизировало бы обучение рассуждению.
\end{itemize}

% ===========================================================
\section*{Упражнения}
\addcontentsline{toc}{section}{Упражнения}
% ===========================================================

\begin{exercise}
\label{ex:rlvr-gradient}
Рассмотрим цель RLVR из~\eqref{eq:rlvr-objective} при $\beta = 0$.
Покажите, что градиент стратегии $J_{\mathrm{RLVR}}(\theta)$ по
$\theta$ может быть записан как
\[
\nabla_\theta J_{\mathrm{RLVR}}(\theta)
  = \E_{x, y}\bigl[\ind\{\textsc{Verify}(y, a^*)\}\,
    \nabla_\theta \log \pi_\theta(y \mid x)\bigr],
\]
начиная с первых принципов с использованием трюка логарифмической производной. Как
это сводится к градиенту REINFORCE\index{REINFORCE} из
Главы~\ref{ch:pg} с $r(y) = \ind\{\textsc{Verify}(y,a^*)\}$?
\end{exercise}

\begin{exercise}
\label{ex:majority-vote-bound}
Докажите Утверждение~\ref{prop:majority-vote}. В частности, покажите, что
для $k$ независимых испытаний Бернулли($p$) при $p > 1/2$ вероятность
того, что произойдёт менее $k/2$ успехов, ограничена сверху
$\exp(-k\,D_{\mathrm{KL}}(1/2\,\|\,1-p))$ с использованием неравенства Хёффдинга
или границ Чернова.
\end{exercise}

\begin{exercise}
\label{ex:grpo-advantage-ch15}
Предположим, что $k = 4$ развёртки для промпта получают вознаграждения $r = (1, 1, 0, 0)$.
Вычислите оценки преимущества GRPO $\hat{A}_i$ из~\eqref{eq:grpo-advantage-ch15}.
Теперь предположим $r = (1, 0, 0, 0)$. Снова вычислите $\hat{A}_i$.
Сравните величину сигнала градиента в двух случаях и
объясните, почему второй сценарий даёт более сильную обратную связь для правильной
развёртки, хотя правильная развёртка только одна.
\end{exercise}

\begin{exercise}
\label{ex:difficulty-filter}
Обучающий набор данных содержит 10\,000 задач. Вы сэмплируете $N = 16$
развёрток на задачу и обнаруживаете, что 2\,000 задач имеют коэффициент решаемости
$\sigma < 0.2$, 5\,000 имеют $\sigma \in [0.2, 0.8]$ и 3\,000
имеют $\sigma > 0.8$.
\begin{enumerate}[label=(\alph*)]
\item Какая доля вычислений «расходуется впустую» на неинформативных задачах
  (с $\sigma < 0.2$ или $\sigma > 0.8$)?
\item Как фильтрация по сложности меняет эффективный размер набора данных,
  доступного для обновлений градиента?
\item Предложите динамический учебный план сложности, который позволяет не отбрасывать
  сложные задачи навсегда.
\end{enumerate}
\end{exercise}

\begin{exercise}
\label{ex:cot-information}
Используя разложение~\eqref{eq:cot-chain-rule}, неформально обоснуйте,
почему (a) chain-of-thought улучшает точность для сложных многошаговых
задач, но (b) по большей части избыточен для задач, которые модель может
решить правильно за один шаг. Что это подразумевает для
требований к \emph{переключаемому} режиму рассуждений?
\end{exercise}

\begin{exercise}
\label{ex:tool-calling-reward}
Модель решает следующую задачу через использование инструментов: получить
численность населения Франции и численность населения Германии, затем вывести их
соотношение. Напишите функцию Python \texttt{verify(trajectory)}, которая принимает
список пар (действие, наблюдение) и возвращает вознаграждение в
$\{0, 1\}$. Ваша функция должна обрабатывать: (a) вызов модели инструментов
в неправильном порядке, (b) отказ модели вызвать требуемый инструмент,
и (c) арифметические ошибки в финальном вычислении.
\end{exercise}

\begin{exercise}
\label{ex:process-vs-outcome}
Рассмотрим многошагового агента, выполняющего задачу из 5 шагов. Вознаграждение за результат
равно $+1$ только при выполнении задачи. Предположим, что агент преуспевает на
финальном шаге, но выбрал неэффективный путь (сделав 3 избыточных API-вызова).
При:
\begin{enumerate}[label=(\alph*)]
\item только вознаграждении за результат, что агент узнаёт об
  избыточных API-вызовах?
\item процессной модели вознаграждения, штрафующей избыточные вызовы на $-0.1$
  каждый, как меняется сигнал градиента?
\end{enumerate}
Обсудите компромисс смещение-дисперсия между (a) и (b).
\end{exercise}

\begin{exercise}
\label{ex:inference-scaling}
Пусть $\text{Acc}_1(p)$~--- точность единичной развёртки, где
каждая развёртка верна с вероятностью $p$. Пусть $\text{Acc}_k(p)$~---
точность голосования большинства по $k$ развёрткам.
\begin{enumerate}[label=(\alph*)]
\item Запишите точную формулу для $\text{Acc}_k(p)$ через
  биномиальное распределение.
\item Для $p = 0.4$ и $k \in \{1, 3, 5, 9, 17, 33\}$ вычислите
  $\text{Acc}_k(p)$. Заметьте, что голосование большинства может \emph{снизить}
  точность при $p < 0.5$; объясните, почему.
\item Найдите пороговое значение $p$, ниже которого голосование большинства вредно.
\end{enumerate}
\end{exercise}

\begin{exercise}
\label{ex:deepseek-recipe}
Рецепт DeepSeek R1 использует два этапа: холодный старт SFT, затем
крупномасштабный RLVR. Предположим, что вы пропускаете этап SFT и применяете RLVR
непосредственно к базовой предобученной модели.
\begin{enumerate}[label=(\alph*)]
\item Какие проблемы могут возникнуть, если базовая модель не надёжно
  производит текст в формате \texttt{<thinking>}...\texttt{</thinking>}?
\item Как награда за формат, описанная в
  Разделе~\ref{subsec:format-rewards}, частично смягчает это?
\item Почему холодный старт SFT всё же предпочтителен, несмотря на
  награду за формат?
\end{enumerate}
\end{exercise}
